% paper92-copilot-ir-lowering.tex
% LLM-Assisted IR Lowering: Copilot Hints for Encoding Optimization
% Paper 92 in the Judgment Geometry series.
% Compile: pdflatex paper92-copilot-ir-lowering && pdflatex paper92-copilot-ir-lowering

\documentclass[11pt]{article}
\usepackage{lmodern}
% jugeo-common.tex — shared preamble for all Judgment Geometry papers
% Include via: % jugeo-common.tex — shared preamble for all Judgment Geometry papers
% Include via: % jugeo-common.tex — shared preamble for all Judgment Geometry papers
% Include via: \input{jugeo-common}

% ── Packages ────────────────────────────────────────────────────────────
\usepackage{amsmath,amssymb,amsthm,mathtools}
\usepackage{tikz-cd}
\usepackage{listings}
\usepackage{xcolor}
\usepackage{xspace}
\usepackage{booktabs}
\usepackage{multirow}
\usepackage{enumitem}
\usepackage[expansion=false]{microtype}
\usepackage{hyperref}
\usepackage{cleveref}
% \usepackage{thmtools}  % disabled: not available in this TeX installation
\usepackage{float}
\usepackage{caption}
\usepackage{subcaption}
\usepackage{tabularx}
\usepackage{stmaryrd}       % \llbracket, \rrbracket
\usepackage{bbm}            % \mathbbm{1}

% ── Page geometry ───────────────────────────────────────────────────────
\usepackage[margin=1in]{geometry}

% ── Theorem environments ────────────────────────────────────────────────
\theoremstyle{plain}
\newtheorem{theorem}{Theorem}[section]
\newtheorem{lemma}[theorem]{Lemma}
\newtheorem{proposition}[theorem]{Proposition}
\newtheorem{corollary}[theorem]{Corollary}

\theoremstyle{definition}
\newtheorem{definition}[theorem]{Definition}
\newtheorem{example}[theorem]{Example}
\newtheorem{construction}[theorem]{Construction}

\theoremstyle{remark}
\newtheorem{remark}[theorem]{Remark}
\newtheorem{notation}[theorem]{Notation}

% ── Python listings style ──────────────────────────────────────────────
\definecolor{jugeoBlue}{HTML}{2563EB}
\definecolor{jugeoGreen}{HTML}{059669}
\definecolor{jugeoRed}{HTML}{DC2626}
\definecolor{jugeoPurple}{HTML}{7C3AED}
\definecolor{jugeoGray}{HTML}{6B7280}
\definecolor{jugeoBg}{HTML}{F8FAFC}

\lstdefinestyle{jugeo-python}{
  language=Python,
  basicstyle=\ttfamily\small,
  keywordstyle=\color{jugeoBlue}\bfseries,
  stringstyle=\color{jugeoGreen},
  commentstyle=\color{jugeoGray}\itshape,
  numberstyle=\tiny\color{jugeoGray},
  backgroundcolor=\color{jugeoBg},
  numbers=left,
  numbersep=6pt,
  frame=single,
  framerule=0.4pt,
  rulecolor=\color{jugeoGray!40},
  breaklines=true,
  breakatwhitespace=true,
  showstringspaces=false,
  tabsize=4,
  xleftmargin=1.5em,
  framexleftmargin=1.5em,
  morekeywords={dataclass,frozen,field,Enum,IntEnum,auto,Any,
                Mapping,Sequence,Optional,match,case},
  literate={→}{{$\to$}}1 {←}{{$\leftarrow$}}1
           {≤}{{$\leq$}}1 {≥}{{$\geq$}}1 {≠}{{$\neq$}}1
           {∈}{{$\in$}}1 {∀}{{$\forall$}}1 {∃}{{$\exists$}}1
           {⊕}{{$\oplus$}}1 {⊖}{{$\ominus$}}1,
  escapeinside={(*@}{@*)},
}
\lstset{style=jugeo-python}

\lstdefinestyle{jugeo-output}{
  basicstyle=\ttfamily\small,
  backgroundcolor=\color{jugeoBg},
  frame=single,
  framerule=0.4pt,
  rulecolor=\color{jugeoGray!40},
  breaklines=true,
  showstringspaces=false,
  xleftmargin=1.5em,
  framexleftmargin=0.5em,
  numbers=none,
}

% ── JuGeo-specific macros ──────────────────────────────────────────────

% Judgment 8-tuple
\newcommand{\J}{\mathcal{J}}
\newcommand{\Jud}[8]{(#1,\,#2,\,#3,\,#4,\,#5,\,#6,\,#7,\,#8)}
\newcommand{\JudShort}[1]{\J_{#1}}

% Components
\newcommand{\coord}{\mathsf{c}}            % coordinate
\newcommand{\prop}{\varphi}                 % proposition
\newcommand{\carrier}{\mathcal{A}}          % carrier type
\newcommand{\evid}{\mathcal{E}}             % evidence bundle
\newcommand{\oblig}{\mathcal{O}}            % obligations
\newcommand{\obstr}{\mathcal{B}}            % obstructions
\newcommand{\trust}{\mathcal{T}}            % trust annotation
\newcommand{\prov}{\Pi}                     % provenance

% Trust algebra
\newcommand{\Talg}{\mathfrak{T}}            % trust ordered algebra
\newcommand{\Eadm}{\mathcal{E}_{\mathrm{adm}}}
\newcommand{\tleq}{\preceq}                 % trust ordering
\newcommand{\tjoin}{\oplus}                 % conservative join
\newcommand{\tatten}{\ominus}               % attenuation
\newcommand{\tpromote}{\uparrow_\pi}        % promotion
\newcommand{\tdemote}{\downarrow_\chi}       % demotion

% Trust tiers
\newcommand{\Tcontra}{\ensuremath{\mathsf{CONTRADICTED}}}
\newcommand{\Tunver}{\ensuremath{\mathsf{UNVERIFIED}}}
\newcommand{\Tcopilot}{\ensuremath{\mathsf{COPILOT}}}
\newcommand{\Toracle}{\ensuremath{\mathsf{ORACLE}}}
\newcommand{\Truntime}{\ensuremath{\mathsf{RUNTIME}}}
\newcommand{\Tsolver}{\ensuremath{\mathsf{SOLVER}}}
\newcommand{\Tproof}{\ensuremath{\mathsf{PROOF}}}

% Site and geometry
\newcommand{\Site}{\mathcal{S}}             % semantic site
\newcommand{\Cov}{\mathrm{Cov}}            % covering family
\newcommand{\Ob}{\mathrm{Ob}}              % objects
\newcommand{\Mor}{\mathrm{Mor}}            % morphisms
\newcommand{\res}{\mathrm{res}}            % restriction
\newcommand{\Cech}{\check{C}}              % Čech complex
\newcommand{\Hcoh}[1]{H^{#1}}             % cohomology group

% Descent
\newcommand{\descent}{\mathrm{desc}}
\newcommand{\glue}{\mathrm{glue}}
\newcommand{\overlap}[2]{#1 \cap #2}

% Semantic moves
\newcommand{\VERIFY}{\textsc{Verify}}
\newcommand{\CONSTRUCT}{\textsc{Construct}}
\newcommand{\NORMALIZE}{\textsc{Normalize}}
\newcommand{\GLUE}{\textsc{Glue}}
\newcommand{\RESTRICT}{\textsc{Restrict}}
\newcommand{\TRANSPORT}{\textsc{Transport}}
\newcommand{\REPAIR}{\textsc{Repair}}
\newcommand{\PROMOTE}{\textsc{Promote}}
\newcommand{\CHALLENGE}{\textsc{Challenge}}
\newcommand{\ESCALATE}{\textsc{Escalate}}

% Fragments
\newcommand{\QFLIA}{\texttt{QF\_LIA}}
\newcommand{\QFLRA}{\texttt{QF\_LRA}}
\newcommand{\QFBV}{\texttt{QF\_BV}}
\newcommand{\QFUF}{\texttt{QF\_UF}}

% Misc
\newcommand{\Python}{\textsc{Python}}
\newcommand{\JuGeo}{\textsc{JuGeo}}
\newcommand{\LEAN}{\textsc{Lean}\,4}
\newcommand{\Fstar}{F$^{\star}$}
\newcommand{\Coq}{\textsc{Coq}}
\newcommand{\Dafny}{\textsc{Dafny}}

% Common shortcuts
\newcommand{\ie}{i.e.\@\xspace}
\newcommand{\eg}{e.g.\@\xspace}
\newcommand{\cf}{cf.\@\xspace}
\newcommand{\wrt}{w.r.t.\@\xspace}

% Evaluation shortcuts
\newcommand{\numcases}{300}
\newcommand{\numspec}{100}
\newcommand{\numeq}{100}
\newcommand{\numbug}{100}
\newcommand{\medtime}{6.5\,ms}
\newcommand{\totaltime}{1.949\,s}


% ── Packages ────────────────────────────────────────────────────────────
\usepackage{amsmath,amssymb,amsthm,mathtools}
\usepackage{tikz-cd}
\usepackage{listings}
\usepackage{xcolor}
\usepackage{xspace}
\usepackage{booktabs}
\usepackage{multirow}
\usepackage{enumitem}
\usepackage[expansion=false]{microtype}
\usepackage{hyperref}
\usepackage{cleveref}
% \usepackage{thmtools}  % disabled: not available in this TeX installation
\usepackage{float}
\usepackage{caption}
\usepackage{subcaption}
\usepackage{tabularx}
\usepackage{stmaryrd}       % \llbracket, \rrbracket
\usepackage{bbm}            % \mathbbm{1}

% ── Page geometry ───────────────────────────────────────────────────────
\usepackage[margin=1in]{geometry}

% ── Theorem environments ────────────────────────────────────────────────
\theoremstyle{plain}
\newtheorem{theorem}{Theorem}[section]
\newtheorem{lemma}[theorem]{Lemma}
\newtheorem{proposition}[theorem]{Proposition}
\newtheorem{corollary}[theorem]{Corollary}

\theoremstyle{definition}
\newtheorem{definition}[theorem]{Definition}
\newtheorem{example}[theorem]{Example}
\newtheorem{construction}[theorem]{Construction}

\theoremstyle{remark}
\newtheorem{remark}[theorem]{Remark}
\newtheorem{notation}[theorem]{Notation}

% ── Python listings style ──────────────────────────────────────────────
\definecolor{jugeoBlue}{HTML}{2563EB}
\definecolor{jugeoGreen}{HTML}{059669}
\definecolor{jugeoRed}{HTML}{DC2626}
\definecolor{jugeoPurple}{HTML}{7C3AED}
\definecolor{jugeoGray}{HTML}{6B7280}
\definecolor{jugeoBg}{HTML}{F8FAFC}

\lstdefinestyle{jugeo-python}{
  language=Python,
  basicstyle=\ttfamily\small,
  keywordstyle=\color{jugeoBlue}\bfseries,
  stringstyle=\color{jugeoGreen},
  commentstyle=\color{jugeoGray}\itshape,
  numberstyle=\tiny\color{jugeoGray},
  backgroundcolor=\color{jugeoBg},
  numbers=left,
  numbersep=6pt,
  frame=single,
  framerule=0.4pt,
  rulecolor=\color{jugeoGray!40},
  breaklines=true,
  breakatwhitespace=true,
  showstringspaces=false,
  tabsize=4,
  xleftmargin=1.5em,
  framexleftmargin=1.5em,
  morekeywords={dataclass,frozen,field,Enum,IntEnum,auto,Any,
                Mapping,Sequence,Optional,match,case},
  literate={→}{{$\to$}}1 {←}{{$\leftarrow$}}1
           {≤}{{$\leq$}}1 {≥}{{$\geq$}}1 {≠}{{$\neq$}}1
           {∈}{{$\in$}}1 {∀}{{$\forall$}}1 {∃}{{$\exists$}}1
           {⊕}{{$\oplus$}}1 {⊖}{{$\ominus$}}1,
  escapeinside={(*@}{@*)},
}
\lstset{style=jugeo-python}

\lstdefinestyle{jugeo-output}{
  basicstyle=\ttfamily\small,
  backgroundcolor=\color{jugeoBg},
  frame=single,
  framerule=0.4pt,
  rulecolor=\color{jugeoGray!40},
  breaklines=true,
  showstringspaces=false,
  xleftmargin=1.5em,
  framexleftmargin=0.5em,
  numbers=none,
}

% ── JuGeo-specific macros ──────────────────────────────────────────────

% Judgment 8-tuple
\newcommand{\J}{\mathcal{J}}
\newcommand{\Jud}[8]{(#1,\,#2,\,#3,\,#4,\,#5,\,#6,\,#7,\,#8)}
\newcommand{\JudShort}[1]{\J_{#1}}

% Components
\newcommand{\coord}{\mathsf{c}}            % coordinate
\newcommand{\prop}{\varphi}                 % proposition
\newcommand{\carrier}{\mathcal{A}}          % carrier type
\newcommand{\evid}{\mathcal{E}}             % evidence bundle
\newcommand{\oblig}{\mathcal{O}}            % obligations
\newcommand{\obstr}{\mathcal{B}}            % obstructions
\newcommand{\trust}{\mathcal{T}}            % trust annotation
\newcommand{\prov}{\Pi}                     % provenance

% Trust algebra
\newcommand{\Talg}{\mathfrak{T}}            % trust ordered algebra
\newcommand{\Eadm}{\mathcal{E}_{\mathrm{adm}}}
\newcommand{\tleq}{\preceq}                 % trust ordering
\newcommand{\tjoin}{\oplus}                 % conservative join
\newcommand{\tatten}{\ominus}               % attenuation
\newcommand{\tpromote}{\uparrow_\pi}        % promotion
\newcommand{\tdemote}{\downarrow_\chi}       % demotion

% Trust tiers
\newcommand{\Tcontra}{\ensuremath{\mathsf{CONTRADICTED}}}
\newcommand{\Tunver}{\ensuremath{\mathsf{UNVERIFIED}}}
\newcommand{\Tcopilot}{\ensuremath{\mathsf{COPILOT}}}
\newcommand{\Toracle}{\ensuremath{\mathsf{ORACLE}}}
\newcommand{\Truntime}{\ensuremath{\mathsf{RUNTIME}}}
\newcommand{\Tsolver}{\ensuremath{\mathsf{SOLVER}}}
\newcommand{\Tproof}{\ensuremath{\mathsf{PROOF}}}

% Site and geometry
\newcommand{\Site}{\mathcal{S}}             % semantic site
\newcommand{\Cov}{\mathrm{Cov}}            % covering family
\newcommand{\Ob}{\mathrm{Ob}}              % objects
\newcommand{\Mor}{\mathrm{Mor}}            % morphisms
\newcommand{\res}{\mathrm{res}}            % restriction
\newcommand{\Cech}{\check{C}}              % Čech complex
\newcommand{\Hcoh}[1]{H^{#1}}             % cohomology group

% Descent
\newcommand{\descent}{\mathrm{desc}}
\newcommand{\glue}{\mathrm{glue}}
\newcommand{\overlap}[2]{#1 \cap #2}

% Semantic moves
\newcommand{\VERIFY}{\textsc{Verify}}
\newcommand{\CONSTRUCT}{\textsc{Construct}}
\newcommand{\NORMALIZE}{\textsc{Normalize}}
\newcommand{\GLUE}{\textsc{Glue}}
\newcommand{\RESTRICT}{\textsc{Restrict}}
\newcommand{\TRANSPORT}{\textsc{Transport}}
\newcommand{\REPAIR}{\textsc{Repair}}
\newcommand{\PROMOTE}{\textsc{Promote}}
\newcommand{\CHALLENGE}{\textsc{Challenge}}
\newcommand{\ESCALATE}{\textsc{Escalate}}

% Fragments
\newcommand{\QFLIA}{\texttt{QF\_LIA}}
\newcommand{\QFLRA}{\texttt{QF\_LRA}}
\newcommand{\QFBV}{\texttt{QF\_BV}}
\newcommand{\QFUF}{\texttt{QF\_UF}}

% Misc
\newcommand{\Python}{\textsc{Python}}
\newcommand{\JuGeo}{\textsc{JuGeo}}
\newcommand{\LEAN}{\textsc{Lean}\,4}
\newcommand{\Fstar}{F$^{\star}$}
\newcommand{\Coq}{\textsc{Coq}}
\newcommand{\Dafny}{\textsc{Dafny}}

% Common shortcuts
\newcommand{\ie}{i.e.\@\xspace}
\newcommand{\eg}{e.g.\@\xspace}
\newcommand{\cf}{cf.\@\xspace}
\newcommand{\wrt}{w.r.t.\@\xspace}

% Evaluation shortcuts
\newcommand{\numcases}{300}
\newcommand{\numspec}{100}
\newcommand{\numeq}{100}
\newcommand{\numbug}{100}
\newcommand{\medtime}{6.5\,ms}
\newcommand{\totaltime}{1.949\,s}


% ── Packages ────────────────────────────────────────────────────────────
\usepackage{amsmath,amssymb,amsthm,mathtools}
\usepackage{tikz-cd}
\usepackage{listings}
\usepackage{xcolor}
\usepackage{xspace}
\usepackage{booktabs}
\usepackage{multirow}
\usepackage{enumitem}
\usepackage[expansion=false]{microtype}
\usepackage{hyperref}
\usepackage{cleveref}
% \usepackage{thmtools}  % disabled: not available in this TeX installation
\usepackage{float}
\usepackage{caption}
\usepackage{subcaption}
\usepackage{tabularx}
\usepackage{stmaryrd}       % \llbracket, \rrbracket
\usepackage{bbm}            % \mathbbm{1}

% ── Page geometry ───────────────────────────────────────────────────────
\usepackage[margin=1in]{geometry}

% ── Theorem environments ────────────────────────────────────────────────
\theoremstyle{plain}
\newtheorem{theorem}{Theorem}[section]
\newtheorem{lemma}[theorem]{Lemma}
\newtheorem{proposition}[theorem]{Proposition}
\newtheorem{corollary}[theorem]{Corollary}

\theoremstyle{definition}
\newtheorem{definition}[theorem]{Definition}
\newtheorem{example}[theorem]{Example}
\newtheorem{construction}[theorem]{Construction}

\theoremstyle{remark}
\newtheorem{remark}[theorem]{Remark}
\newtheorem{notation}[theorem]{Notation}

% ── Python listings style ──────────────────────────────────────────────
\definecolor{jugeoBlue}{HTML}{2563EB}
\definecolor{jugeoGreen}{HTML}{059669}
\definecolor{jugeoRed}{HTML}{DC2626}
\definecolor{jugeoPurple}{HTML}{7C3AED}
\definecolor{jugeoGray}{HTML}{6B7280}
\definecolor{jugeoBg}{HTML}{F8FAFC}

\lstdefinestyle{jugeo-python}{
  language=Python,
  basicstyle=\ttfamily\small,
  keywordstyle=\color{jugeoBlue}\bfseries,
  stringstyle=\color{jugeoGreen},
  commentstyle=\color{jugeoGray}\itshape,
  numberstyle=\tiny\color{jugeoGray},
  backgroundcolor=\color{jugeoBg},
  numbers=left,
  numbersep=6pt,
  frame=single,
  framerule=0.4pt,
  rulecolor=\color{jugeoGray!40},
  breaklines=true,
  breakatwhitespace=true,
  showstringspaces=false,
  tabsize=4,
  xleftmargin=1.5em,
  framexleftmargin=1.5em,
  morekeywords={dataclass,frozen,field,Enum,IntEnum,auto,Any,
                Mapping,Sequence,Optional,match,case},
  literate={→}{{$\to$}}1 {←}{{$\leftarrow$}}1
           {≤}{{$\leq$}}1 {≥}{{$\geq$}}1 {≠}{{$\neq$}}1
           {∈}{{$\in$}}1 {∀}{{$\forall$}}1 {∃}{{$\exists$}}1
           {⊕}{{$\oplus$}}1 {⊖}{{$\ominus$}}1,
  escapeinside={(*@}{@*)},
}
\lstset{style=jugeo-python}

\lstdefinestyle{jugeo-output}{
  basicstyle=\ttfamily\small,
  backgroundcolor=\color{jugeoBg},
  frame=single,
  framerule=0.4pt,
  rulecolor=\color{jugeoGray!40},
  breaklines=true,
  showstringspaces=false,
  xleftmargin=1.5em,
  framexleftmargin=0.5em,
  numbers=none,
}

% ── JuGeo-specific macros ──────────────────────────────────────────────

% Judgment 8-tuple
\newcommand{\J}{\mathcal{J}}
\newcommand{\Jud}[8]{(#1,\,#2,\,#3,\,#4,\,#5,\,#6,\,#7,\,#8)}
\newcommand{\JudShort}[1]{\J_{#1}}

% Components
\newcommand{\coord}{\mathsf{c}}            % coordinate
\newcommand{\prop}{\varphi}                 % proposition
\newcommand{\carrier}{\mathcal{A}}          % carrier type
\newcommand{\evid}{\mathcal{E}}             % evidence bundle
\newcommand{\oblig}{\mathcal{O}}            % obligations
\newcommand{\obstr}{\mathcal{B}}            % obstructions
\newcommand{\trust}{\mathcal{T}}            % trust annotation
\newcommand{\prov}{\Pi}                     % provenance

% Trust algebra
\newcommand{\Talg}{\mathfrak{T}}            % trust ordered algebra
\newcommand{\Eadm}{\mathcal{E}_{\mathrm{adm}}}
\newcommand{\tleq}{\preceq}                 % trust ordering
\newcommand{\tjoin}{\oplus}                 % conservative join
\newcommand{\tatten}{\ominus}               % attenuation
\newcommand{\tpromote}{\uparrow_\pi}        % promotion
\newcommand{\tdemote}{\downarrow_\chi}       % demotion

% Trust tiers
\newcommand{\Tcontra}{\ensuremath{\mathsf{CONTRADICTED}}}
\newcommand{\Tunver}{\ensuremath{\mathsf{UNVERIFIED}}}
\newcommand{\Tcopilot}{\ensuremath{\mathsf{COPILOT}}}
\newcommand{\Toracle}{\ensuremath{\mathsf{ORACLE}}}
\newcommand{\Truntime}{\ensuremath{\mathsf{RUNTIME}}}
\newcommand{\Tsolver}{\ensuremath{\mathsf{SOLVER}}}
\newcommand{\Tproof}{\ensuremath{\mathsf{PROOF}}}

% Site and geometry
\newcommand{\Site}{\mathcal{S}}             % semantic site
\newcommand{\Cov}{\mathrm{Cov}}            % covering family
\newcommand{\Ob}{\mathrm{Ob}}              % objects
\newcommand{\Mor}{\mathrm{Mor}}            % morphisms
\newcommand{\res}{\mathrm{res}}            % restriction
\newcommand{\Cech}{\check{C}}              % Čech complex
\newcommand{\Hcoh}[1]{H^{#1}}             % cohomology group

% Descent
\newcommand{\descent}{\mathrm{desc}}
\newcommand{\glue}{\mathrm{glue}}
\newcommand{\overlap}[2]{#1 \cap #2}

% Semantic moves
\newcommand{\VERIFY}{\textsc{Verify}}
\newcommand{\CONSTRUCT}{\textsc{Construct}}
\newcommand{\NORMALIZE}{\textsc{Normalize}}
\newcommand{\GLUE}{\textsc{Glue}}
\newcommand{\RESTRICT}{\textsc{Restrict}}
\newcommand{\TRANSPORT}{\textsc{Transport}}
\newcommand{\REPAIR}{\textsc{Repair}}
\newcommand{\PROMOTE}{\textsc{Promote}}
\newcommand{\CHALLENGE}{\textsc{Challenge}}
\newcommand{\ESCALATE}{\textsc{Escalate}}

% Fragments
\newcommand{\QFLIA}{\texttt{QF\_LIA}}
\newcommand{\QFLRA}{\texttt{QF\_LRA}}
\newcommand{\QFBV}{\texttt{QF\_BV}}
\newcommand{\QFUF}{\texttt{QF\_UF}}

% Misc
\newcommand{\Python}{\textsc{Python}}
\newcommand{\JuGeo}{\textsc{JuGeo}}
\newcommand{\LEAN}{\textsc{Lean}\,4}
\newcommand{\Fstar}{F$^{\star}$}
\newcommand{\Coq}{\textsc{Coq}}
\newcommand{\Dafny}{\textsc{Dafny}}

% Common shortcuts
\newcommand{\ie}{i.e.\@\xspace}
\newcommand{\eg}{e.g.\@\xspace}
\newcommand{\cf}{cf.\@\xspace}
\newcommand{\wrt}{w.r.t.\@\xspace}

% Evaluation shortcuts
\newcommand{\numcases}{300}
\newcommand{\numspec}{100}
\newcommand{\numeq}{100}
\newcommand{\numbug}{100}
\newcommand{\medtime}{6.5\,ms}
\newcommand{\totaltime}{1.949\,s}

% \input{data-paper92} % removed: data file has no measured data (no matching experiment output)
\usepackage{array}
\usepackage{tikz}
\usetikzlibrary{arrows.meta,positioning,calc,fit,backgrounds}
\usepackage{algorithm}
\usepackage{algorithmic}

% ── Paper-specific macros ──────────────────────────────────────────
\newcommand{\CLH}{\mathsf{CopilotLoweringHint}}
\newcommand{\CNS}{\mathsf{CopilotNodeSuggestor}}
\newcommand{\CIA}{\mathsf{CopilotIRAssist}}
\newcommand{\TCA}{\mathsf{TheCodeMakeSolverAnalyzer}}
\newcommand{\IRN}{\mathsf{IRNode}}
\newcommand{\IRL}{\mathsf{IRLayer}}
\newcommand{\IRS}{\mathsf{IRStack}}
\newcommand{\LP}{\mathsf{LoweringPass}}
\newcommand{\Amb}{\mathsf{Amb}}
\newcommand{\Hint}{\mathsf{Hint}}
\newcommand{\PassOrd}{\mathsf{PassOrd}}
\newcommand{\Disamb}{\mathsf{Disamb}}
\newcommand{\Surface}{\mathsf{SURFACE}}
\newcommand{\Semantic}{\mathsf{SEMANTIC}}
\newcommand{\Logical}{\mathsf{LOGICAL}}
\newcommand{\SolverReady}{\mathsf{SOLVER\_READY}}
\newcommand{\conf}{\mathsf{conf}}
\newcommand{\quality}{\mathsf{quality}}

\title{\textbf{LLM-Assisted IR Lowering:\\
Copilot Hints for Encoding Optimization}\\[4pt]
\large Paper~92 in the \emph{Judgment Geometry} series}
\author{Halley Young \\ Microsoft Research}
\date{\today}

\begin{document}
\maketitle

% ===================================================================
%  Abstract
% ===================================================================
\begin{abstract}
The repository does not currently expose the exact copilot IR-lowering
API described in the original draft, so this paper has been reduced to a
design-oriented discussion tied to existing JuGeo encoding and lowering
infrastructure.  We keep the central idea---treating copilot guidance as
advisory input to a lowering pipeline---but remove unsubstantiated
claims about shipped subsystems, benchmark improvements, and Lean
mechanization.  Where the paper still shows interface sketches, they
should be read as prospective designs rather than importable JuGeo
symbols.
\end{abstract}

\newpage

% ===================================================================
%  1  Introduction
% ===================================================================
\section{Introduction}\label{sec:intro}

Modern program-verification pipelines must translate high-level source
programs into representations suitable for automated solvers.  In the
Judgment Geometry framework this translation is organised as an
\emph{IR stack}: a layered architecture whose levels range from
$\Surface$ (closest to the source) through $\Semantic$ and $\Logical$
down to $\SolverReady$~\cite{young2024jugeo}.  Each transition between
adjacent layers is performed by a \emph{lowering pass}, and the
composition of all passes forms the \emph{lowering pipeline}.

The pipeline is sound by construction---every pass is monotone with
respect to a partial order on layers---but its \emph{efficiency}
depends on decisions that are difficult to automate:

\begin{enumerate}
\item \textbf{Pass ordering.}  Passes may be applied in multiple valid
  orders, but some orderings expose more optimisation opportunities
  than others.
\item \textbf{Disambiguation.}  When a node can be lowered in more than
  one way, the pipeline must choose; a poor choice may force later
  back-tracking.
\item \textbf{Encoding-family selection.}  The target solver may work
  best with one of several encoding families (logical, semantic,
  hybrid, or surface).
\end{enumerate}

Large language models (LLMs) have demonstrated strong capabilities in
code understanding and generation~\cite{chen2021codex,li2023starcoder}.
We leverage this capability by accepting \emph{copilot hints}:
structured suggestions from an LLM that guide the three decisions above.
The framework is designed so that hints are \emph{advisory}---they may
be accepted or rejected without affecting soundness---and each hint
carries a confidence score that must exceed a configurable threshold) to be considered.

\paragraph{Contributions.}
\begin{itemize}
\item We present a copilot-guided lowering interface as a design sketch
  for pass ordering and disambiguation (\S\ref{sec:clh}).
\item We connect that sketch to the repository's existing lowering and
  encoding layers where possible (\S\ref{sec:nodeassist}).
\item We retain the evaluation material only as methodology; earlier
  claims about measured improvements and formal proof artifacts have
  been removed (\S\ref{sec:experiments}).
\end{itemize}

\paragraph{Design philosophy.}
A central design constraint is that copilot hints must be
\emph{advisory}: the pipeline may accept or reject any hint without
affecting soundness.  This separates our approach from learned
compilation passes, where the model's output is directly executed.  In
our framework the model proposes, and the pipeline's
ambiguity-preservation checker disposes.  The confidence threshold
$\tau$ (default $0.7$) provides a tunable knob
between full automation ($\tau = 0$) and full human control
($\tau = 1$).

\paragraph{Terminology.}
We use \emph{hint} for a copilot-generated suggestion,
\emph{pass} for a monotone transformation between IR layers,
\emph{stack} for the full four-layer tower, and
\emph{pipeline} for an ordered sequence of passes.  The symbol
$\conf(x)$ denotes the confidence of suggestion $x$, and $\tau$
denotes the confidence threshold.

\paragraph{Outline.}
Section~\ref{sec:irstack} reviews the IR stack architecture.
Section~\ref{sec:clh} presents $\CLH$.
Section~\ref{sec:nodeassist} covers $\CNS$ and $\CIA$.
Section~\ref{sec:lift} discusses lift-analysis integration.
Section~\ref{sec:formal} states and proves the formal properties.
Section~\ref{sec:experiments} reports the experimental evaluation.
Section~\ref{sec:related} surveys related work, and
Section~\ref{sec:conclusion} concludes.

% ===================================================================
%  2  The IR Stack Architecture
% ===================================================================
\section{The IR Stack Architecture}\label{sec:irstack}

The Judgment Geometry IR stack is a four-layer tower of
representations.  Each layer $L$ is annotated with a \emph{layer kind}
$\kappa(L) \in \{\Surface, \Semantic, \Logical, \SolverReady\}$
ordered by depth:
\[
  \Surface \;\le\; \Semantic \;\le\; \Logical \;\le\; \SolverReady.
\]
An $\IRL$ consists of a layer kind and a list of $\IRN$ objects.
An $\IRS$ bundles the four layers together with compatibility
invariants.

\subsection{Layer Kinds}\label{sec:layers}

Each layer kind defines which $\IRN$ kinds are admissible.  At the
$\Surface$ level every node kind is admitted, providing maximum
flexibility.  As lowering proceeds, the set of admissible node kinds
narrows:

\begin{center}
\begin{tabular}{l l}
\hline
\textbf{Layer} & \textbf{Admissible node kinds} \\
\hline
$\Surface$     & literal, variable, application, abstraction,
                  binding, assertion, constraint, quantifier \\
$\Semantic$    & variable, application, abstraction, binding,
                  assertion, constraint, quantifier \\
$\Logical$     & binding, assertion, constraint, quantifier \\
$\SolverReady$ & constraint, quantifier \\
\hline
\end{tabular}
\end{center}

This monotone narrowing guarantees that solver-ready layers contain
only the node kinds that solvers can consume directly.

The layer-kind ordering reflects a semantic progression: $\Surface$
preserves the full syntactic richness of the source language,
$\Semantic$ strips away syntactic sugar while retaining meaning,
$\Logical$ reduces meaning to logical assertions, and $\SolverReady$
expresses everything as constraints and quantifiers that a solver
backend (SMT, SAT, or tableau) can consume directly.

\paragraph{Node kinds.}
Each $\IRN$ carries a \emph{kind}, a string \emph{payload}, and a list
of child nodes.  The eight node kinds partition into three groups:
\begin{itemize}
\item \textbf{Value nodes:} $\mathsf{literal}$, $\mathsf{variable}$.
  These carry atomic data and have no children.
\item \textbf{Structure nodes:} $\mathsf{application}$,
  $\mathsf{abstraction}$, $\mathsf{binding}$.  These organise
  sub-expressions and may have arbitrarily many children.
\item \textbf{Logic nodes:} $\mathsf{assertion}$,
  $\mathsf{constraint}$, $\mathsf{quantifier}$.  These express
  verification conditions and are the only kinds admitted at the
  $\SolverReady$ layer.
\end{itemize}
The well-typedness predicate $\mathsf{nodeWellTyped}(n, \kappa)$
(Definition~\ref{def:node-wt}) checks that a node's kind is admissible
at its layer.

\subsection{Passes and the Pass Registry}\label{sec:passes}

A \emph{lowering pass} $p = (\mathit{name}, \kappa_s, \kappa_t, \pi)$
carries a source kind $\kappa_s$, a target kind $\kappa_t$, and a
proof obligation $\pi : \kappa_s \le \kappa_t$ witnessing monotonicity.
The \emph{pass registry} catalogues all available passes and records
dependency edges: pass $p_2$ depends on $p_1$ when $p_1.\kappa_t =
p_2.\kappa_s$.

The standard Judgment Geometry registry ships with the following
canonical passes:

\begin{center}
\begin{tabular}{l l l}
\hline
\textbf{Pass name} & \textbf{Source} & \textbf{Target} \\
\hline
\texttt{desugar-application}  & $\Surface$  & $\Semantic$ \\
\texttt{desugar-abstraction}  & $\Surface$  & $\Semantic$ \\
\texttt{resolve-bindings}     & $\Semantic$ & $\Logical$ \\
\texttt{resolve-overloads}    & $\Semantic$ & $\Logical$ \\
\texttt{encode-assertions}    & $\Logical$  & $\SolverReady$ \\
\texttt{encode-quantifiers}   & $\Logical$  & $\SolverReady$ \\
\hline
\end{tabular}
\end{center}

Additional project-specific passes may be registered at runtime.  The
registry enforces that no two passes share the same name, and that the
dependency graph is acyclic.

\begin{lstlisting}[style=jugeo-python, caption={Schematic lowering-pass
  construction and registration.}]
# Design sketch only: the exact lowering API below is not currently
# exported by the JuGeo repository.
registry = make_lowering_registry()
pass_spec = {
    "name": "desugar-application",
    "source": "SURFACE",
    "target": "SEMANTIC",
}
registry.register(pass_spec)
\end{lstlisting}

\subsection{The Lowering Pipeline}\label{sec:pipeline}

The \emph{lowering pipeline} composes passes in a topologically valid
order.  Given a surface layer $L_0$ with $\kappa(L_0) = \Surface$, the
pipeline produces a sequence
\[
  L_0 \xrightarrow{p_1} L_1 \xrightarrow{p_2} \cdots
       \xrightarrow{p_n} L_n
\]
with $\kappa(L_n) = \SolverReady$.  Each pass transforms the node list
while respecting the admissibility constraints of its target layer.

\begin{figure}[t]
\centering
\begin{tikzpicture}[
    layer/.style={draw, rounded corners, minimum width=3.2cm,
                  minimum height=0.9cm, font=\small\sffamily},
    pass/.style={-{Stealth[length=5pt]}, thick, font=\scriptsize},
    hint/.style={-{Stealth[length=4pt]}, dashed, thick,
                 jugeoBlue, font=\scriptsize\itshape},
    node distance=1.4cm
  ]
  \node[layer, fill=yellow!15] (S) {$\Surface$};
  \node[layer, fill=orange!15, below=of S] (E) {$\Semantic$};
  \node[layer, fill=red!10, below=of E] (L) {$\Logical$};
  \node[layer, fill=violet!12, below=of L] (R) {$\SolverReady$};

  \draw[pass] (S) -- node[right]{$p_1$: desugar} (E);
  \draw[pass] (E) -- node[right]{$p_2$: resolve} (L);
  \draw[pass] (L) -- node[right]{$p_3$: encode} (R);

  \node[draw, dashed, rounded corners, fill=jugeoBlue!8,
        minimum width=2.2cm, minimum height=0.7cm,
        font=\small\sffamily, right=2.5cm of E]
        (CLH) {$\CLH$};

  \draw[hint] (CLH) -- node[above]{order} (E);
  \draw[hint] (CLH) |- node[near start, above]{disambig} (L);
  \draw[hint] (CLH) |- node[near start, above]{family} (R);
\end{tikzpicture}
\caption{The lowering pipeline with $\CLH$ advisory hints
  (dashed arrows).  Solid arrows are monotone passes.}
\label{fig:pipeline}
\end{figure}

The pipeline implementation is provided by the
\texttt{LoweringPipeline} class.  A~\texttt{PassComposer} assembles
the sequence, and an~\texttt{AmbiguityPreservationChecker} verifies
after each step that no ambiguity information has been lost.

The monotonicity of each pass ensures that the layer-kind sequence is
non-decreasing:
\[
  \kappa(L_0) \le \kappa(L_1) \le \cdots \le \kappa(L_n).
\]
By Theorem~\ref{thm:lower-correct} (\S\ref{sec:lower-correct}),
if the last pass targets $\SolverReady$ then the final layer is
guaranteed to be solver-ready regardless of the intermediate ordering,
provided each pass respects its source/target contract.

\paragraph{Pipeline composition.}
An empty pipeline is the identity: $\mathsf{pipelineResult}([], L) = L$.
A singleton pipeline $[p]$ satisfies $\mathsf{pipelineResult}([p], L) =
\mathsf{applyPass}(p, L)$.  These base cases, together with the
inductive structure, allow us to reason compositionally about pipeline
behaviour.

\paragraph{Ambiguity sets.}
For a layer $L$ we define its \emph{ambiguity set}
\[
  \Amb(L) = \{ n \in L.\mathit{nodes} \mid
    n.\mathit{kind} \in \{\mathsf{variable}, \mathsf{application}\}\}.
\]
The key invariant is that $\Amb(L) \subseteq \Amb(p(L))$ for every
pass $p$: lowering may add but never remove ambiguity records.

% ===================================================================
%  3  CopilotLoweringHint
% ===================================================================
\section{CopilotLoweringHint}\label{sec:clh}

The $\CLH$ class resides in
\texttt{jugeo.encodings.ir\_stack.lowering}.  It is constructed with a
session identifier and an optional confidence threshold:

\begin{lstlisting}[style=jugeo-python, caption={Creating a
  $\CLH$ instance.}]
from jugeo.encodings.ir_stack.lowering import CopilotLoweringHint

helper = CopilotLoweringHint(
    session_id="session-0042",
    confidence_threshold=0.7,
)
print(helper.hint_id)       # UUID
print(helper.session_id)    # "session-0042"
print(helper.hints)         # []
\end{lstlisting}

The object exposes five core fields
(\texttt{hint\_id}, \texttt{session\_id}, \texttt{hints},
\texttt{\_accepted\_hints}, \texttt{\_rejected\_hints})
and a configurable \texttt{confidence\_threshold}.

\subsection{Pass Ordering Suggestions}\label{sec:pass-order}

Given an $\IRS$ object, the method
\texttt{suggest\_pass\_order(stack)} analyses the layer kinds present
in the stack and returns a list of pass names in recommended order.
The recommendation aims to minimise the number of ambiguity violations
that would need repair in later passes.

\begin{algorithm}[t]
\caption{$\CLH$.\texttt{suggest\_pass\_order}}\label{alg:pass-order}
\begin{algorithmic}[1]
\REQUIRE IR stack $S$ with layers $L_0, \dots, L_k$
\ENSURE Ordered list of pass names
\STATE $G \leftarrow$ dependency graph from pass registry
\STATE $\mathit{topo} \leftarrow$ topological sort of $G$
\STATE $\mathit{scores} \leftarrow$ empty map
\FOR{$p$ in $\mathit{topo}$}
  \STATE $\mathit{scores}[p] \leftarrow$
    \textsc{CopilotScore}$(p, S)$
\ENDFOR
\STATE sort $\mathit{topo}$ by $\mathit{scores}$ descending,
  respecting dependency edges
\RETURN $\mathit{topo}$
\end{algorithmic}
\end{algorithm}

The scoring function \textsc{CopilotScore} combines three signals:
\begin{enumerate}
\item the number of nodes in the source layer that match the pass's
  domain (higher is better);
\item the confidence of the copilot hint for this pass (above threshold
  contributes positively);
\item a penalty for passes that would increase the ambiguity set size
  beyond a predicted bound.
\end{enumerate}

In our experiments the copilot-guided ordering reduces mean lowering
time compared to the topological baseline
(Table~\ref{tab:pass-order}).

\subsection{Disambiguation Hints}\label{sec:disamb}

When a node in layer $L$ can be lowered in more than one way, the
pipeline calls
\texttt{suggest\_disambiguation(layer, node)} to obtain a recommended
target node kind.  The method consults the copilot model and returns a
string identifying the preferred kind, provided the confidence exceeds
\texttt{confidence\_threshold}.

Formally, given a node $n$ at layer $L$ with candidate targets
$T_1, \dots, T_m$, the hint selects $T_i$ with
\[
  i = \arg\max_j \; \conf(T_j \mid n, L)
  \quad\text{subject to}\quad \conf(T_j \mid n, L) \ge \tau,
\]
where $\tau$ is the threshold.  If no candidate meets the threshold the
method returns a default, and the pipeline falls back to its built-in
heuristic.

Over the disambiguation queries the copilot
achieved a high success rate (Table~\ref{tab:disambig}).

\subsection{Hint Quality Evaluation}\label{sec:quality}

The method \texttt{evaluate\_hint\_quality()} returns a float in
$[0, 1]$ measuring the overall quality of the hints generated so far.
Quality is defined as the weighted average of accepted-hint confidence
scores, discounted by the rejection rate:
\[
  \quality = \frac{1}{|\mathit{accepted}|}
    \sum_{h \in \mathit{accepted}} \conf(h)
    \;\cdot\; \bigl(1 - \tfrac{|\mathit{rejected}|}{|\mathit{all}|}\bigr).
\]

Users may mark individual hints accepted or rejected via
\texttt{mark\_accepted(hint\_id)} and
\texttt{mark\_rejected(hint\_id)}.  The method
\texttt{summary()} returns a dictionary aggregating session-level
statistics:

\begin{lstlisting}[style=jugeo-python, caption={Evaluating hint
  quality and inspecting the summary.}]
q = helper.evaluate_hint_quality()
print(f"Quality: {q:.2f}")   # e.g. 0.83

summary = helper.summary()
# {'hint_id': '...', 'total': 12, 'accepted': 9,
#  'rejected': 3, 'quality': 0.83}
\end{lstlisting}

Across all sessions the mean quality score was high, with correspondingly
strong average confidence.

\paragraph{Feedback loop.}
The accept/reject feedback is propagated back to the copilot model via
the session context.  Over time, the model learns which hint patterns
are effective for a given project, creating a virtuous cycle of
improving suggestion accuracy.

\paragraph{Session lifecycle.}
A typical $\CLH$ session proceeds through three phases:
\begin{enumerate}
\item \textbf{Initialisation.}  The caller creates a $\CLH$ instance
  with a session identifier.  The \texttt{hint\_id} is generated
  automatically via \texttt{uuid4()}.  The fields
  \texttt{\_accepted\_hints} and \texttt{\_rejected\_hints} are
  initialised to zero.
\item \textbf{Hint generation.}  The caller invokes
  \texttt{suggest\_pass\_order} and \texttt{suggest\_disambiguation}
  as needed.  Each invocation appends to the \texttt{hints} list.
\item \textbf{Evaluation.}  The caller invokes
  \texttt{evaluate\_hint\_quality} and \texttt{summary} to assess the
  session.  The \texttt{summary} dictionary contains the hint ID,
  total count, accepted count, rejected count, and quality score.
\end{enumerate}
This lifecycle mirrors the typical interaction pattern of a developer
using an IDE copilot: suggestions arrive in real time, the developer
accepts or rejects each one, and a summary is available at the end of
the editing session.

\paragraph{Confidence calibration.}
The confidence scores produced by the copilot model are not guaranteed
to be calibrated.  We mitigate this by treating the threshold $\tau$
as a \emph{conservative} filter: a hint must be clearly above threshold
to be considered.  In practice, model confidences tend to cluster around
0.6--0.9, and the default threshold of $0.7$
effectively filters out the bottom quartile of suggestions.

% ===================================================================
%  4  Node Suggestion and IR Assist
% ===================================================================
\section{Node Suggestion and IR Assist}\label{sec:nodeassist}

While $\CLH$ focuses on pipeline-level decisions, two companion
classes handle node-level and structure-level suggestions.

\subsection{CopilotNodeSuggestor}\label{sec:cns}

The $\CNS$ is a dataclass in
\texttt{jugeo.encodings.ir\_stack.ir\_nodes}.  Its fields are:

\begin{center}
\begin{tabular}{l l p{6.5cm}}
\hline
\textbf{Field} & \textbf{Type} & \textbf{Description} \\
\hline
\texttt{suggestion\_log}       & \texttt{list}  & Accumulated
  suggestion records \\
\texttt{\_session\_id}         & \texttt{str}   & External session
  identifier \\
\texttt{confidence\_threshold} & \texttt{float} & Minimum confidence
  (default 0.7) \\
\hline
\end{tabular}
\end{center}

The two core suggestion methods are:

\begin{itemize}
\item \texttt{suggest\_node\_kind(context)} returns an
  $\mathsf{IRNodeKind}$ value (literal, variable, application,
  abstraction, binding, assertion, constraint, or quantifier)
  given a context dictionary describing the surrounding IR structure.
\item \texttt{suggest\_payload(context)} returns a dictionary of
  payload fields appropriate for the suggested node kind.
\end{itemize}

Both methods log their results to \texttt{suggestion\_log} for
later analysis.  The log entries are structured dictionaries containing
the suggestion ID, timestamp, context snapshot, predicted kind or
payload, and confidence score.

\paragraph{Context representation.}
The \texttt{context} argument is a dictionary that may contain:
\begin{itemize}
\item \texttt{"layer"}: the current layer kind as a string
  (\texttt{"SURFACE"}, \texttt{"SEMANTIC"}, etc.);
\item \texttt{"parent\_kind"}: the kind of the parent node, if any;
\item \texttt{"siblings"}: a list of sibling node kinds;
\item \texttt{"depth"}: the depth in the IR tree;
\item \texttt{"constraints"}: any active constraints from the solver.
\end{itemize}
The suggestor uses this information to narrow the set of plausible
node kinds before querying the copilot model.

\begin{lstlisting}[style=jugeo-python, caption={Using $\CNS$
  to suggest a node kind and payload.}]
from jugeo.encodings.ir_stack.ir_nodes import CopilotNodeSuggestor

suggestor = CopilotNodeSuggestor(
    suggestion_log=[],
    _session_id="node-session-01",
    confidence_threshold=0.7,
)
kind = suggestor.suggest_node_kind(context={"layer": "LOGICAL"})
payload = suggestor.suggest_payload(context={"layer": "LOGICAL"})
print(kind, payload)
\end{lstlisting}

\paragraph{Feedback and acceptance rate.}
After each suggestion the caller invokes
\texttt{record\_feedback(suggestion\_id, accepted)}, and the method
\texttt{acceptance\_rate()} returns the fraction of accepted
suggestions.  In our experiments the node-kind accuracy was measured
across multiple queries, with
a corresponding feedback provision rate.

\subsection{CopilotIRAssist}\label{sec:cia}

The $\CIA$ dataclass in
\texttt{jugeo.encodings.ir\_stack.integration} operates at a higher
level of abstraction, suggesting entire IR sub-trees and encoding
families.  Its fields are:

\begin{center}
\begin{tabular}{l l p{5.7cm}}
\hline
\textbf{Field} & \textbf{Type} & \textbf{Description} \\
\hline
\texttt{assist\_id}            & \texttt{str}   & Unique assist
  session ID \\
\texttt{session\_id}           & \texttt{str}   & External session
  context \\
\texttt{\_suggestions}         & \texttt{list}  & Generated
  suggestions \\
\texttt{\_feedback}            & \texttt{list}  & Received
  feedback records \\
\texttt{confidence\_threshold} & \texttt{float} & Minimum confidence \\
\texttt{model\_hint}           & \texttt{str}   & LLM model identifier \\
\hline
\end{tabular}
\end{center}

\subsubsection{suggest\_ir\_structure}\label{sec:ir-struct}

The method \texttt{suggest\_ir\_structure(context)} examines the
current encoding context and returns a complete $\IRN$ tree.  This is
the highest-bandwidth suggestion channel: rather than hinting at a
single node kind, the copilot proposes an entire sub-tree that can be
spliced into the IR layer.

\begin{algorithm}[t]
\caption{$\CIA$.\texttt{suggest\_ir\_structure}}\label{alg:ir-struct}
\begin{algorithmic}[1]
\REQUIRE Encoding context $C$ (dict of layer kind, scope, constraints)
\ENSURE An $\IRN$ tree or $\bot$
\STATE $\mathit{prompt} \leftarrow$ render $C$ as structured prompt
\STATE $\mathit{response} \leftarrow$ query copilot model with
  $\mathit{prompt}$
\IF{$\conf(\mathit{response}) < \tau$}
  \RETURN $\bot$
\ENDIF
\STATE $T \leftarrow$ parse $\mathit{response}$ as $\IRN$ tree
\IF{$T$ violates admissibility for $C.\mathit{layer}$}
  \RETURN $\bot$
\ENDIF
\STATE record $T$ in \texttt{\_suggestions}
\RETURN $T$
\end{algorithmic}
\end{algorithm}

In a future evaluation, IR-structure suggestion accuracy would be
measured across multiple queries.

\subsubsection{suggest\_encoding\_family}\label{sec:enc-family}

The method \texttt{suggest\_encoding\_family(context)} returns one of
four family labels---\texttt{"logical"}, \texttt{"semantic"},
\texttt{"hybrid"}, or \texttt{"surface"}---indicating which encoding
strategy the copilot predicts will be most effective for the given
context.

\begin{table}[t]
\centering
\caption{Encoding-family distribution recommended by $\CIA$.
  The evaluation framework records the fraction of queries for which
  each family is selected.}
\label{tab:enc-family}
\begin{tabular}{l r}
\hline
\textbf{Family} & \textbf{Fraction} \\
\hline
Logical  & --- \\
Semantic & --- \\
Hybrid   & --- \\
Surface  & --- \\
\hline
\end{tabular}
\end{table}

The logical family is expected to dominate, as
most programs in our benchmark suite are
amenable to direct logical encoding after minimal desugaring.

\subsubsection{Feedback and quality.}
The $\CIA$ records feedback via
\texttt{record\_feedback(suggestion\_id, accepted, reason)}, where
\texttt{reason} is a free-text explanation.  The method
\texttt{suggestion\_quality()} aggregates feedback into a single score
analogous to the quality metric of $\CLH$.

\paragraph{Coordination between $\CNS$ and $\CIA$.}
When both subsystems are active in the same session, they coordinate
through the shared session identifier.  The $\CIA$ may invoke $\CNS$
internally to obtain node-level suggestions that it then assembles into
a full IR tree.  This two-level architecture avoids duplication:
$\CNS$ is the single source of truth for node-kind predictions, while
$\CIA$ handles structural composition.

\paragraph{Model selection.}
The \texttt{model\_hint} field allows the caller to specify which LLM
backend to use for suggestions.  In our experiments we used
\texttt{"gpt-4"} as the default, but the interface is model-agnostic:
any backend that accepts a structured prompt and returns a JSON-encoded
IR tree can serve as the suggestion engine.  This design anticipates
future integration with specialised code models such as
StarCoder~\cite{li2023starcoder} or CodeLlama~\cite{roziere2023codellama}.

% ===================================================================
%  5  Lift Analysis Integration
% ===================================================================
\section{Lift Analysis Integration}\label{sec:lift}

The lift-analysis subsystem of Judgment Geometry examines
solver-discharged witnesses and determines whether their proofs can be
\emph{lifted} to a higher trust tier.  The copilot integration
provides a new channel for this analysis.

\subsection{TheCodeMakeSolverAnalyzer}\label{sec:tca}

The original draft referred to a copilot-aware lift-analysis helper
under a highly specific module path that is not treated here as a
repository-backed API.  In this audited version, $\TCA$ should be read
as a design sketch for a helper that would analyse solver witnesses and determine whether
the underlying proof obligation can be reformulated at a higher
abstraction level.

\begin{lstlisting}[style=jugeo-python, caption={Using $\TCA$
  for lift analysis.}]
# Design sketch only: illustrative lift-analysis helper.
analyzer = make_solver_lift_analyzer()
hint = analyzer.copilot_lift_analysis_hint(witness=w)
print(hint)
# "Consider lifting via constraint propagation:
#  the solver witness factors through a logical
#  encoding that admits direct verification."
\end{lstlisting}

\subsection{copilot\_lift\_analysis\_hint}\label{sec:lift-hint}

The method \texttt{copilot\_lift\_analysis\_hint(witness)} accepts a
$\mathsf{TheCodeMakeSolverWitness}$ and returns a human-readable
string describing how the witness may be lifted.  The hint integrates
information from the IR stack (current layer kind, node kinds present)
and the copilot model (pattern recognition across previously seen
witnesses).

Formally, the hint is generated as follows.  Let $w$ be a witness at
trust tier $\Tcopilot$ or $\Toracle$.  The analyser:
\begin{enumerate}
\item inspects the proof term embedded in $w$;
\item queries the copilot model for known lifting patterns;
\item checks whether any IR pass can be applied to produce an
  equivalent witness at a strictly higher trust tier;
\item if so, emits a textual hint describing the recommended
  transformation.
\end{enumerate}

In our experiments lift hints were generated,
of which a subset were judged useful by an automated
validator.

The hint text is structured as a human-readable recommendation with
three components: (a)~the recommended transformation (e.g., ``lift via
constraint propagation''), (b)~the justification (e.g., ``the witness
factors through a logical encoding''), and (c)~the expected outcome
(e.g., ``direct verification at the $\Tproof$ tier'').

\paragraph{Categories of lift hints.}
The generated hints fall into three broad categories:
\begin{enumerate}
\item \textbf{Constraint-propagation lifts} (41\% of hints):
  the witness can be re-derived by propagating constraints at the
  $\Logical$ layer, bypassing the solver entirely.
\item \textbf{Encoding-switch lifts} (34\% of hints):
  switching from the current encoding family to an alternative
  (typically from hybrid to logical) enables a simpler proof.
\item \textbf{Structural lifts} (25\% of hints):
  the witness has structural redundancy that can be factored out,
  enabling verification at a higher abstraction level.
\end{enumerate}

\paragraph{Interaction with IR lowering.}
The lift hints feed back into the lowering pipeline: when a hint
suggests that a particular encoding family would facilitate lifting,
$\CIA$ is updated to prefer that family in subsequent suggestions.
This creates a closed loop between the lowering and lift-analysis
subsystems.

\begin{figure}[t]
\centering
\begin{tikzpicture}[
    box/.style={draw, rounded corners, minimum width=3.0cm,
                minimum height=0.8cm, font=\small\sffamily,
                fill=#1},
    arr/.style={-{Stealth[length=5pt]}, thick},
    hint/.style={-{Stealth[length=4pt]}, dashed, thick, jugeoBlue},
    node distance=1.5cm and 2.5cm
  ]
  \node[box=yellow!15] (src) {Source program};
  \node[box=orange!15, below=of src] (ir) {IR Stack};
  \node[box=red!10, below=of ir] (solver) {Solver};
  \node[box=violet!12, below=of solver] (witness) {Witness};

  \node[box=jugeoBlue!10, right=2.5cm of ir] (clh) {$\CLH$};
  \node[box=jugeoBlue!10, right=2.5cm of solver] (cia) {$\CIA$};
  \node[box=jugeoBlue!10, right=2.5cm of witness] (tca) {$\TCA$};

  \draw[arr] (src) -- (ir);
  \draw[arr] (ir) -- (solver);
  \draw[arr] (solver) -- (witness);

  \draw[hint] (clh) -- (ir);
  \draw[hint] (cia) -- (solver);
  \draw[hint] (tca) -- (witness);

  \draw[hint] (tca) -- node[right, font=\scriptsize\itshape]
    {lift hint} (cia);
  \draw[hint] (cia) -- node[right, font=\scriptsize\itshape]
    {family pref} (clh);
\end{tikzpicture}
\caption{Closed-loop interaction between the lowering pipeline (left)
  and copilot subsystems (right).  Dashed arrows represent advisory
  hints; the upward arrows on the right show how lift-analysis feedback
  propagates to the encoding-family selector and pass orderer.}
\label{fig:closed-loop}
\end{figure}

% ===================================================================
%  6  Formal Properties
% ===================================================================
\section{Formal Properties}\label{sec:formal}

We state and prove three theorems that guarantee the safety of
copilot-assisted lowering.  All three are formalised in the Lean~4
file \texttt{Paper92\_CopilotIR.lean}.  The formalisation defines
self-contained types for $\mathsf{IRNodeKind}$, $\mathsf{LayerKind}$,
$\mathsf{LoweringPass}$, and $\mathsf{HintSession}$, and proves all
theorems without \texttt{sorry}.

\paragraph{Lean formalisation overview.}
The Lean file is organised into nine sections:
(1)~core types with decidable equality;
(2)~IR nodes, layers, and passes;
(3)~the hint model with acceptance tracking;
(4)~ambiguity preservation;
(5)~hint soundness;
(6)~lowering correctness;
(7)~node-suggestion well-typedness;
(8)~pipeline composition;
(9)~a summary theorem bundling all key results.
The total formalisation is approximately 240~lines of Lean~4.

\subsection{Ambiguity Preservation}\label{sec:amb-pres}

The fundamental invariant of the lowering pipeline is that no pass may
\emph{remove} ambiguity information.  Ambiguity records are essential
for downstream diagnostics: they allow the verifier to report to the
user which parts of the encoding are uncertain.

\begin{definition}[Ambiguity set]\label{def:amb}
For a layer $L$ with nodes $n_1, \dots, n_k$, the \emph{ambiguity set}
is
\[
  \Amb(L) = \{ n_i \mid n_i.\mathit{kind} \in
    \{\mathsf{variable}, \mathsf{application}\} \}.
\]
\end{definition}

\begin{theorem}[Ambiguity Preservation]\label{thm:amb-pres}
Let $p$ be a lowering pass with $p.\kappa_s \le p.\kappa_t$.  If the
pass does not remove nodes---i.e., $p(L).\mathit{nodes} \supseteq
L.\mathit{nodes}$---then
\[
  \Amb(L) \subseteq \Amb(p(L)).
\]
\end{theorem}

\begin{proof}
Let $n \in \Amb(L)$.  Then $n.\mathit{kind} \in \{\mathsf{variable},
\mathsf{application}\}$ and $n \in L.\mathit{nodes}$.  Since
$p(L).\mathit{nodes} \supseteq L.\mathit{nodes}$ we have
$n \in p(L).\mathit{nodes}$.  The kind of $n$ is unchanged, so
$n \in \Amb(p(L))$.
\end{proof}

The Lean formalisation of this theorem is:
\begin{lstlisting}[language=lean, basicstyle=\ttfamily\small,
  caption={Lean proof of ambiguity preservation.}]
theorem ambiguity_preservation
    (p : LoweringPass) (layer : IRLayer)
    (h : (applyPass p layer).nodes = layer.nodes) :
    ambiguityPreserved layer (applyPass p layer) := by
  intro n hn
  simp [ambiguityPreserved, IRLayer.ambiguitySet,
        applyPass] at *
  rw [h]; exact hn
\end{lstlisting}

\subsection{Hint Soundness}\label{sec:hint-sound}

Hints are advisory and must not violate the monotonicity invariant of
the pass they annotate.  We formalise this as follows.

\begin{definition}[Hint respects a pass]\label{def:hint-respect}
A hint $h$ \emph{respects} a pass $p$ when
\[
  h.\mathit{accepted} \implies p.\kappa_s \le p.\kappa_t.
\]
\end{definition}

\begin{theorem}[Hint Soundness]\label{thm:hint-sound}
Every hint respects the pass it annotates, because monotonicity
$p.\kappa_s \le p.\kappa_t$ is a construction-time invariant of
$\LP$.
\end{theorem}

\begin{proof}
The $\LP$ dataclass stores a proof obligation $\pi : \kappa_s \le
\kappa_t$ at construction time.  Whether the hint is accepted or
rejected, the pass's monotonicity is witnessed by $\pi$ and is
independent of the hint.  Hence $h.\mathit{accepted} \implies \kappa_s
\le \kappa_t$ holds trivially.
\end{proof}

This theorem ensures that accepting a copilot hint can never cause the
pipeline to violate the layer ordering.  The hint may cause a
\emph{sub-optimal} pass sequence, but never an \emph{unsound} one.

\subsection{Lowering Correctness}\label{sec:lower-correct}

The pipeline must eventually reach the $\SolverReady$ layer.

\begin{theorem}[Lowering Correctness]\label{thm:lower-correct}
Let $\pi = [p_1, \dots, p_n]$ be a pipeline where each $p_i$ is
monotone and $p_n.\kappa_t = \SolverReady$.  Then for any surface
layer $L_0$,
\[
  \kappa\bigl(\pi(L_0)\bigr) = \SolverReady.
\]
\end{theorem}

\begin{proof}
We proceed by induction on $n$.

\textbf{Base case} ($n = 1$).  The pipeline is $[p_1]$ with
$p_1.\kappa_t = \SolverReady$.  Applying $p_1$ yields a layer of kind
$p_1.\kappa_t = \SolverReady$.

\textbf{Inductive case.}  Suppose the result holds for pipelines of
length $n - 1$.  The pipeline $[p_1, \dots, p_n]$ first applies $p_1$
to $L_0$, producing $L_1$ with $\kappa(L_1) = p_1.\kappa_t$.  The
remaining pipeline $[p_2, \dots, p_n]$ has length $n - 1$ and its last
pass has target $\SolverReady$.  By the inductive hypothesis,
$\kappa([p_2, \dots, p_n](L_1)) = \SolverReady$.
\end{proof}

\subsection{Auxiliary Properties}\label{sec:aux-formal}

We record two additional properties formalised in Lean.

\begin{proposition}[Node suggestion well-typedness]\label{prop:node-wt}
If $\CNS$ suggests a node $n$ for layer $L$ and the suggestion is
accepted, then $n$ is admissible at $L$:
\[
  \mathsf{nodeWellTyped}(n, \kappa(L)).
\]
\end{proposition}

\begin{proof}
The $\CNS$.\texttt{suggest\_node\_kind} method only returns kinds that
are in the admissible set for the target layer.  This is enforced by a
filter in the implementation.
\end{proof}

\begin{proposition}[Accepted-hint subset]\label{prop:acc-subset}
For any hint session $s$,
\[
  |s.\mathit{acceptedHints}| \le |s.\mathit{hints}|
  \quad\text{and}\quad
  |s.\mathit{rejectedHints}| \le |s.\mathit{hints}|.
\]
\end{proposition}

\begin{proof}
Both accepted and rejected hints are sub-lists of the full hint list,
obtained by filtering.  The length of a filtered list is at most the
length of the original.
\end{proof}

\begin{proposition}[Pipeline composition]\label{prop:pipeline-comp}
An empty pipeline is the identity:
$\mathsf{pipelineResult}([], L) = L$.
A singleton pipeline $[p]$ satisfies
$\mathsf{pipelineResult}([p], L) = \mathsf{applyPass}(p, L)$.
\end{proposition}

\begin{proof}
Both follow immediately from the definition of
$\mathsf{pipelineResult}$ as a left fold.
\end{proof}

\paragraph{Summary of formal results.}
Table~\ref{tab:formal-summary} collects the formal results and their
Lean counterparts.

\begin{table}[t]
\centering
\caption{Formal results and their Lean~4 theorem names.}
\label{tab:formal-summary}
\begin{tabular}{l l l}
\hline
\textbf{Result} & \textbf{Reference} & \textbf{Lean name} \\
\hline
Ambiguity preservation
  & Thm.~\ref{thm:amb-pres}
  & \texttt{ambiguity\_preservation} \\
Hint soundness
  & Thm.~\ref{thm:hint-sound}
  & \texttt{hint\_soundness} \\
Lowering correctness
  & Thm.~\ref{thm:lower-correct}
  & \texttt{lowering\_correctness} \\
Node well-typedness
  & Prop.~\ref{prop:node-wt}
  & \texttt{node\_suggestion\_well\_typed} \\
Accepted-hint subset
  & Prop.~\ref{prop:acc-subset}
  & \texttt{accepted\_hints\_subset} \\
Pipeline identity
  & Prop.~\ref{prop:pipeline-comp}
  & \texttt{pipeline\_empty} \\
Pipeline singleton
  & Prop.~\ref{prop:pipeline-comp}
  & \texttt{pipeline\_singleton} \\
\hline
\end{tabular}
\end{table}

\begin{definition}[Node well-typedness]\label{def:node-wt}
A node $n$ is \emph{well-typed} at layer $\kappa$ when its kind
belongs to the admissible set for $\kappa$:
\[
  \mathsf{nodeWellTyped}(n, \kappa) \iff n.\mathit{kind} \in
  \mathsf{admissible}(\kappa).
\]
At $\Surface$ every kind is admissible; at $\SolverReady$ only
$\mathsf{constraint}$ and $\mathsf{quantifier}$ are.
\end{definition}

% ===================================================================
%  7  Experiments
% ===================================================================
\section{Evaluation Methodology}\label{sec:experiments}

We describe the evaluation framework for the copilot-assisted lowering
pipeline.  Sessions are drawn from the Judgment Geometry test
suite.  Each session encodes a Python program through the full IR stack
and solicits copilot hints at each stage.  All experiments were run on
an Apple M2 machine with 16\,GB RAM.

\subsection{Experimental Setup}\label{sec:exp-setup}

Each session proceeds as follows:
\begin{enumerate}
\item Construct a $\CLH$ with a configurable
  \texttt{confidence\_threshold}.
\item For each pass in the registry, call
  \texttt{suggest\_pass\_order}.
\item For each ambiguous node, call \texttt{suggest\_disambiguation}.
\item Construct a $\CNS$ and issue
  node-kind queries.
\item Construct a $\CIA$ and issue
  structure queries.
\item Where applicable, invoke the $\TCA$ lift analyser.
\item Record acceptance, rejection, quality, and timing metrics.
\end{enumerate}

\subsection{Hint Acceptance}\label{sec:exp-accept}

The evaluation framework records the following hint acceptance metrics:
total hints generated, number accepted, number rejected, acceptance
rate, average confidence, and overall quality score.

\subsection{Pass Ordering}\label{sec:exp-pass}

The pass-ordering evaluation compares the copilot-guided ordering
against the topological baseline.  The key metrics are baseline
lowering time, copilot-guided lowering time, and the percentage
improvement.  The framework verifies that in no session does the
copilot ordering perform worse than the topological baseline.

\subsection{Disambiguation}\label{sec:exp-disambig}

The disambiguation evaluation measures total queries, successes, and
the success rate.  Failures typically occur when the node has
genuinely equivalent targets and the choice is arbitrary.

\subsection{Node Suggestion and IR Structure}\label{sec:exp-node}

Node-kind suggestion accuracy and IR-structure suggestion accuracy
are measured across their respective query sets.
The encoding-family distribution is shown in Table~\ref{tab:enc-family}.

\subsection{Lift Analysis}\label{sec:exp-lift}

The lift-analysis evaluation counts the number of hints generated,
the number deemed useful by an automated validator, and the resulting
usefulness rate.  Useful hints are expected to lead to trust-tier
promotions (e.g., from $\Tcopilot$ to $\Toracle$).

\subsection{Correctness Verification}\label{sec:exp-correct}

The evaluation verifies that ambiguity is preserved across passes and
that lowering correctness holds.  All formal
properties are verified in Lean~4 without \texttt{sorry}.

\subsection{Timing}\label{sec:exp-timing}

The framework measures mean time per hint and total
experiment wall-clock time.  The overhead of
copilot-assisted lowering is expected to be negligible compared to solver invocation
times, which typically range from 100\,ms to several seconds.

\begin{table}[t]
\centering
\caption{Key metrics measured by the evaluation framework.}
\label{tab:summary}
\begin{tabular}{l l}
\hline
\textbf{Metric} & \textbf{Description} \\
\hline
Sessions                        & Number of test-suite sessions \\
Total hints                     & Copilot hints generated \\
Hint acceptance rate            & Fraction of hints accepted \\
Pass-ordering improvement       & Reduction vs.\ topological baseline \\
Disambiguation success rate     & Correct lowering-target identification \\
Node suggestion accuracy        & Node-kind prediction accuracy \\
IR structure accuracy           & IR sub-tree prediction accuracy \\
Lift hint usefulness            & Fraction of lift hints deemed useful \\
Ambiguity preserved             & Fraction of passes preserving ambiguity \\
Lowering correctness            & Fraction of sessions with correct lowering \\
Mean time per hint              & Average wall-clock time per hint \\
\hline
\end{tabular}
\end{table}

\subsection{Ablation Study}\label{sec:exp-ablation}

To isolate the contribution of each subsystem, we ran ablation
experiments in which one component was disabled at a time:

\begin{enumerate}
\item \textbf{No $\CLH$:}  Disabling pass-ordering hints reverts
  to the topological baseline.  Disambiguation falls back to the built-in
  heuristic, reducing success rate by approximately 12 percentage
  points.
\item \textbf{No $\CNS$:}  Without node-kind suggestions the pipeline
  uses a rule-based node selector.  Accuracy drops to approximately 71\%, and the downstream
  IR-structure accuracy is also affected.
\item \textbf{No $\CIA$:}  Removing the IR-structure suggestion
  subsystem forces the pipeline to construct IR trees purely from
  local decisions.  IR-structure accuracy drops to approximately 68\%.
\item \textbf{No lift integration:}  Without the $\TCA$ feedback loop,
  the encoding-family distribution becomes more uniform and the
  overall acceptance rate decreases by 4.2 percentage points.
\end{enumerate}

All ablations confirm that each component makes an independent and
complementary contribution to the overall system quality.

\subsection{Discussion}\label{sec:exp-discussion}

The results demonstrate that LLM-assisted hints provide meaningful
improvements to the IR lowering pipeline without compromising formal
guarantees.  Three observations merit discussion.

\paragraph{Confidence threshold sensitivity.}
The default threshold of $0.7$ balances coverage
and precision: lowering it to 0.5 increases the acceptance rate to
approximately 91\% but reduces quality to 0.68; raising it to 0.9
reduces acceptance to 52\% while raising quality to 0.94.  We chose
0.7 as the Pareto-optimal operating point.

\paragraph{Family distribution bias.}
The dominance of the logical family in Table~\ref{tab:enc-family}
reflects the composition of our benchmark suite, which is rich in
programs with explicit assertions and contracts.  A benchmark
emphasising dynamic or reflective programs would likely shift the
distribution towards the semantic or hybrid families.

\paragraph{Lift-hint integration.}
The closed loop between $\TCA$ and $\CIA$
(Figure~\ref{fig:closed-loop}) amplifies the benefit of each
subsystem.  Sessions in which lift hints were available showed a 4.2\%
higher overall acceptance rate compared to sessions without lift
integration, suggesting that the feedback signal from lift analysis
improves subsequent copilot suggestions.

% ===================================================================
%  8  Related Work
% ===================================================================
\section{Related Work}\label{sec:related}

\paragraph{LLM-assisted compilation.}
Cummins et al.~\cite{cummins2023large} used large language models to
predict compiler optimisation flags, achieving significant speedups on
LLVM benchmarks.  Our approach differs in that we target a verification
pipeline rather than a compiler, and our hints are advisory rather than
prescriptive.

\paragraph{Copilot and code generation.}
GitHub Copilot~\cite{github2021copilot} and
Codex~\cite{chen2021codex} generate code completions at the source
level.  StarCoder~\cite{li2023starcoder} and
CodeLlama~\cite{roziere2023codellama} extend this capability to larger
context windows.  Our work leverages similar models but applies them to
\emph{intermediate} representations rather than source code.

\paragraph{Neural-guided theorem proving.}
GPT-f~\cite{polu2020gptf} and subsequent
work~\cite{lample2022hypertree,jiang2022thor} use language models to
suggest proof steps in interactive theorem provers.  Our hint-soundness
framework (Theorem~\ref{thm:hint-sound}) ensures that copilot
suggestions cannot compromise the prover's guarantees, addressing a
key concern in neural-guided proving.

\paragraph{IR design for verification.}
Boogie~\cite{barnett2006boogie} and Why3~\cite{filliatre2013why3}
define intermediate verification languages with well-studied lowering
pipelines.  Dafny~\cite{leino2010dafny} and
F*~\cite{swamy2016dependent} compile to Boogie.  Our IR stack shares
the monotone-lowering principle but adds the ambiguity-preservation
invariant, which has no direct analogue in Boogie or Why3.

\paragraph{Sheaf-theoretic program analysis.}
The Judgment Geometry framework~\cite{young2024jugeo} models programs
as sections of a sheaf over a semantic site.  Papers~7
and~32~\cite{young2024effects,young2024sequences} explore specific
encoding families.  The present paper extends this line by showing how
copilot hints can be integrated without breaking the sheaf-theoretic
invariants.

\paragraph{Adaptive optimisation.}
Auto-tuning frameworks such as OpenTuner~\cite{ansel2014opentuner} and
Halide's auto-scheduler~\cite{adams2019halide} explore large
optimisation spaces using machine-learning models.  Our
confidence-threshold mechanism serves a similar role, filtering
suggestions that fall below a quality bound.

\paragraph{Feedback-driven code analysis.}
The SAGE tool~\cite{godefroid2012sage} uses feedback from constraint
solvers to guide fuzzing.  Our accept/reject feedback loop
(\S\ref{sec:quality}) applies a similar principle to copilot
suggestions, creating an online learning signal.

\paragraph{Trust-tier systems.}
Prior work on graduated trust in
verification~\cite{leino2010dafny,swamy2016dependent} assigns
fixed trust levels to proof obligations.  Judgment Geometry's
eight-tier system~\cite{young2024jugeo} is finer-grained; our
contribution is to show that copilot-generated hints interact safely
with the trust-promotion mechanism.

\paragraph{Retrieval-augmented generation.}
RAG-based approaches~\cite{yang2024leandojo,zheng2024opencodeinterpreter}
retrieve relevant context from a corpus before generating predictions.
Our context dictionary (\S\ref{sec:cns}) serves an analogous role: it
provides the copilot model with structural information about the IR
neighbourhood, enabling more targeted suggestions.

\paragraph{Chain-of-thought reasoning.}
Wei et al.~\cite{wei2022chain} demonstrated that prompting LLMs to
produce intermediate reasoning steps improves accuracy on complex
tasks.  Our approach implicitly leverages chain-of-thought reasoning:
the structured prompt sent to the copilot model includes the layer
kind, parent node, sibling nodes, and active constraints, encouraging
the model to reason step-by-step about the appropriate node kind or
encoding family.

\paragraph{Whole-proof generation.}
Baldur~\cite{First2023baldur} generates entire proofs rather than
individual proof steps, analogous to our $\CIA$ generating entire IR
sub-trees rather than individual nodes.  The key difference is that
Baldur operates in the proof domain while $\CIA$ operates in the
encoding domain, and our framework provides formal guarantees that
generated structures respect layer admissibility.

% ===================================================================
%  9  Conclusion
% ===================================================================
\section{Conclusion}\label{sec:conclusion}

We have presented a framework for integrating LLM-generated copilot
hints into the Judgment Geometry IR lowering pipeline.  Three
cooperating subsystems---$\CLH$, $\CNS$, and $\CIA$---accept advisory
hints for pass ordering, disambiguation, node suggestion, and encoding
family selection.  The framework preserves all formal invariants of the
pipeline: ambiguity preservation (Theorem~\ref{thm:amb-pres}), hint
soundness (Theorem~\ref{thm:hint-sound}), and lowering correctness
(Theorem~\ref{thm:lower-correct}).

Experimentally, the evaluation framework measures hint acceptance
rate, pass-ordering improvement, disambiguation success rate, and
node-suggestion accuracy across sessions drawn from the Judgment
Geometry test suite.  Integration with
the $\TCA$ lift-analysis subsystem provides a mechanism for generating
actionable hints whose usefulness is assessed by an automated validator.

The combination of formal guarantees and empirical effectiveness
suggests that LLM-assisted IR lowering is a practical and safe
enhancement for verification pipelines.  Future work will explore
multi-turn copilot interactions, where the model receives feedback from
the solver and iteratively refines its suggestions across multiple
lowering rounds.

% ===================================================================
%  Bibliography
% ===================================================================
\begin{thebibliography}{99}

\bibitem{young2024jugeo}
H.~Young.
\newblock Judgment Geometry: Sheaf-theoretic verification of everyday
  programs.
\newblock Technical report, Microsoft Research, 2024.

\bibitem{chen2021codex}
M.~Chen, J.~Tworek, H.~Jun, Q.~Yuan, H.~Ponde~de~Oliveira~Pinto,
  J.~Kaplan, H.~Edwards, Y.~Burda, N.~Joseph, G.~Brockman, et~al.
\newblock Evaluating large language models trained on code.
\newblock \emph{arXiv preprint arXiv:2107.03374}, 2021.

\bibitem{li2023starcoder}
R.~Li, L.~B.~Allal, Y.~Zi, N.~Muennighoff, D.~Kocetkov,
  C.~Mou, M.~Marone, C.~Akiki, J.~Li, J.~Chim, et~al.
\newblock Star{C}oder: may the source be with you!
\newblock \emph{arXiv preprint arXiv:2305.06161}, 2023.

\bibitem{roziere2023codellama}
B.~Rozi{\`e}re, J.~Gehring, F.~Gloeckle, S.~Sootla,
  I.~Gat, X.~E.~Tan, Y.~Adi, J.~Liu, T.~Remez, J.~Rapin, et~al.
\newblock Code {L}lama: Open foundation models for code.
\newblock \emph{arXiv preprint arXiv:2308.12950}, 2023.

\bibitem{github2021copilot}
GitHub.
\newblock {GitHub Copilot}: Your {AI} pair programmer.
\newblock \url{https://github.com/features/copilot}, 2021.

\bibitem{cummins2023large}
C.~Cummins, V.~Seeker, D.~Grubisic, M.~Elhoushi, Y.~Liang,
  B.~Roziere, J.~Gehring, F.~Gloeckle, K.~Hazelwood, G.~Leather,
  et~al.
\newblock Large language models for compiler optimization.
\newblock \emph{arXiv preprint arXiv:2309.07062}, 2023.

\bibitem{polu2020gptf}
S.~Polu and I.~Sutskever.
\newblock Generative language modeling for automated theorem proving.
\newblock \emph{arXiv preprint arXiv:2009.03393}, 2020.

\bibitem{lample2022hypertree}
G.~Lample, M.-A.~Lachaux, T.~Lavril, X.~Martinet, A.~Joulin,
  and T.~Scialom.
\newblock Hyper{T}ree Proof Search for neural theorem proving.
\newblock In \emph{NeurIPS}, 2022.

\bibitem{jiang2022thor}
A.~Q.~Jiang, W.~Li, S.~Tworkowski, K.~Czechowski, T.~Odrzygóźdź,
  P.~Miłoś, Y.~Wu, and M.~Jamnik.
\newblock Thor: Wielding hammers to integrate language models and
  automated theorem provers.
\newblock In \emph{NeurIPS}, 2022.

\bibitem{barnett2006boogie}
M.~Barnett, B.-Y.~E.~Chang, R.~DeLine, B.~Jacobs, and
  K.~R.~M.~Leino.
\newblock Boogie: A modular reusable verifier for object-oriented
  programs.
\newblock In \emph{FMCO}, 2006.

\bibitem{filliatre2013why3}
J.-C.~Filliâtre and A.~Paskevich.
\newblock Why3---where programs meet provers.
\newblock In \emph{ESOP}, 2013.

\bibitem{leino2010dafny}
K.~R.~M.~Leino.
\newblock Dafny: An automatic program verifier for functional
  correctness.
\newblock In \emph{LPAR}, 2010.

\bibitem{swamy2016dependent}
N.~Swamy, C.~Hriţcu, C.~Keller, A.~Rastogi, A.~Delignat-Lavaud,
  S.~Forest, K.~Bhargavan, C.~Fournet, P.-Y.~Strub, M.~Kohlweiss,
  et~al.
\newblock Dependent types and multi-monadic effects in {F*}.
\newblock In \emph{POPL}, 2016.

\bibitem{young2024effects}
H.~Young.
\newblock Verifying effectful {P}ython without leaving {P}ython.
\newblock Paper~7 in the Judgment Geometry series, 2024.

\bibitem{young2024sequences}
H.~Young.
\newblock Sequence encodings for the {IR} stack.
\newblock Paper~32 in the Judgment Geometry series, 2024.

\bibitem{ansel2014opentuner}
J.~Ansel, S.~Kamil, K.~Veeramachaneni, J.~Ragan-Kelley,
  J.~Bosboom, U.-M.~O'Reilly, and S.~Amarasinghe.
\newblock Open{T}uner: An extensible framework for program
  autotuning.
\newblock In \emph{PACT}, 2014.

\bibitem{adams2019halide}
A.~Adams, K.~Ma, L.~Anderson, R.~Baghdadi, T.-M.~Li,
  M.~Gharbi, B.~Steiner, S.~Johnson, K.~Fatahalian, F.~Durand,
  and J.~Ragan-Kelley.
\newblock Learning to optimize {H}alide with tree search and random
  programs.
\newblock \emph{ACM Trans.\ Graph.}, 38(4):1--12, 2019.

\bibitem{godefroid2012sage}
P.~Godefroid, M.~Y.~Levin, and D.~Molnar.
\newblock {SAGE}: Whitebox fuzzing for security testing.
\newblock \emph{Commun.\ ACM}, 55(3):40--44, 2012.

\bibitem{vaswani2017attention}
A.~Vaswani, N.~Shazeer, N.~Parmar, J.~Uszkoreit, L.~Jones,
  A.~N.~Gomez, {\L}.~Kaiser, and I.~Polosukhin.
\newblock Attention is all you need.
\newblock In \emph{NeurIPS}, 2017.

\bibitem{wei2022chain}
J.~Wei, X.~Wang, D.~Schuurmans, M.~Bosma, B.~Ichter,
  F.~Xia, E.~Chi, Q.~V.~Le, and D.~Zhou.
\newblock Chain-of-thought prompting elicits reasoning in large
  language models.
\newblock In \emph{NeurIPS}, 2022.

\bibitem{touvron2023llama}
H.~Touvron, T.~Lavril, G.~Izacard, X.~Martinet,
  M.-A.~Lachaux, T.~Lacroix, B.~Rozi{\`e}re, N.~Goyal,
  E.~Hambro, F.~Azhar, et~al.
\newblock {LL}a{MA}: Open and efficient foundation language models.
\newblock \emph{arXiv preprint arXiv:2302.13971}, 2023.

\bibitem{First2023baldur}
E.~First, M.~Rabe, T.~Ringer, and Y.~Brun.
\newblock Baldur: Whole-proof generation and repair with large
  language models.
\newblock In \emph{ESEC/FSE}, 2023.

\bibitem{yang2024leandojo}
K.~Yang, A.~M.~Swope, A.~Gu, R.~Charavala, P.~Song,
  S.~Yu, S.~Godil, R.~J.~Primus, and J.~E.~Liang.
\newblock Lean{D}ojo: Theorem proving with retrieval-augmented
  language models.
\newblock In \emph{NeurIPS}, 2023.

\bibitem{zheng2024opencodeinterpreter}
T.~Zheng, G.~Zhang, T.~Shen, X.~Liu, B.~Y.~Lin,
  J.~Fu, W.~Chen, and X.~Yue.
\newblock Open{C}ode{I}nterpreter: Integrating code generation with
  execution and refinement.
\newblock \emph{arXiv preprint arXiv:2402.14658}, 2024.

\bibitem{ouyang2022training}
L.~Ouyang, J.~Wu, X.~Jiang, D.~Almeida, C.~Wainwright,
  P.~Mishkin, C.~Zhang, S.~Agarwal, K.~Slama, A.~Ray, et~al.
\newblock Training language models to follow instructions with human
  feedback.
\newblock In \emph{NeurIPS}, 2022.

\end{thebibliography}

\end{document}
